\section{投资调度求解与完整时序验证}

\subsection{成本、总碳与硬边界}

投资、网络选择与设备模式包含整数变量，运行状态和能流为连续变量，实际装配后形成MILP。成本目标写为
\begin{equation}
F_C=C_{\mathrm{cap}}^{\mathrm{ann}}+C_{\mathrm{energy}}
+C_{\mathrm{om}}+C_{\mathrm{carbon}}+C_{\mathrm{pen}},
\label{eq:cost-objective}
\end{equation}
其中资本年化参数、能源与运维价格、假设碳价和惩罚定义均来自数据字典。已启用的火用偏好成本归入惩罚项并单列。$F_E$为冷热服务与共享电力一致口径的系统总碳。可靠性、舒适和资源边界不与成本目标交换；同一费用或排放项只进入一个明确的组件台账。

\subsection{词典序端点与\texorpdfstring{$\varepsilon$}{epsilon}约束}

端点与碳帽扫描参考多目标约束法的组织原则\supcite{5}，具体容差和时限身份由本项目实际实现与运行记录确定。

模型提供两阶段词典序接口：对主目标$f_1$先求解，得到可行值$v_1$和上下界；再保留原约束，并增加
\begin{equation}
f_1\leq v_1+\delta_{\mathrm{abs}}+\delta_{\mathrm{rel}}|v_1|,
\label{eq:lexicographic}
\end{equation}
随后最小化次目标$f_2$。成本端点采用$(f_1,f_2)=(F_C,F_E)$，碳端点采用$(F_E,F_C)$；两个端点使用独立模型实例，记录两阶段状态、gap、容差和求解时间。第一阶段时间受限时，第二阶段仅在已获得可行值的容差边界内优化，不能称为已证明的全局词典序最优。

$\varepsilon$约束在统一可靠性与服务边界下求解
\begin{equation}
\min F_C,\qquad F_E\leq\varepsilon_E.
\label{eq:epsilon-constraint}
\end{equation}
两个扫描入口的目标分别记录：\path{experiments.py}的专用扫描将碳价设为零；\path{run_competition_frontier.py}保留情景配置的总目标，包括碳价成本及已启用的偏好项。因此赛事扫描结果不能标为纯粹的无碳价成本前沿。碳帽原值和数值容差后的有效值同时保存。检查每档状态、总碳及数值可行性后，对实际采样可行点剔除重复和支配解；有限采样集合不能写成完整精确Pareto前沿，也不声称实现了完整AUGMECON。

上述接口需由对应实验记录证明。单目标运行且未设置碳帽的manifest，不构成已完成$\varepsilon$扫描或词典序端点证明；供热原型的多目标运行与后述公开数学例题复现分别验收。

\subsection{代表期规划与投资冻结}

代表期保留气象、得热、舒适和服务状态的同源组合，记录实际日到代表日的映射、小时权重及覆盖审计。独立典型日难以表达跨期储热联系，代表期聚合须处理状态的时间联系\supcite{4}。本稿不声称实现该文的季节储能状态分解；状态边界单独给出，不能拼接不同年度输入以补齐所需小时，也不能将加权代表小时表述为完整连续运行。

先在代表期确定投资方案，保存容量、技术、网络、储热、接入和初末边界等记录。全年及供暖季运行沿用同一投资指纹和服务合同，新增或退役投资变量冻结；两种时序的范围和状态边界各自声明。完整时序由实际输入时间戳构成，运行层重新调度但不重新选择投资。

求解通过SciPy 1.15.3的\path{scipy.optimize.milp}接口调用HiGHS\supcite{6}。日志记录状态、原始与对偶界、相对gap、时限和数值审计；终止与可行性容差按接口及实际求解器选项解释\supcite{7}，不假定包装接口支持全部原生选项。时限内可行解保留其近似身份；通过运行审计仅说明所给投资在该时序内可运行，不说明全年度投资已联合优化到全局最优。

\subsection{独立数值与来源审核}

矩阵审计检查约束、变量上下界、整数可行性、目标重算及碳帽。建筑轨迹由独立RC复算核对采样温度、边界、冷热互斥和交付量。储热、逐源年度地源、资本一次计入和冷热总碳分别形成审计台账。

独立复算须使用可追溯的序列、配置和输入锚点，留存源文件与代码epoch指纹；只审查已保存矩阵的报告，不能宣称独立重建了全部设备、网络方程。来源审核另检查字段依据和变换身份。数值通过、来源通过和外部观测验证是不同证据，不能相互代替。

\clearpage
\section{验证证据、适用范围与对照}

\subsection{验证范围}

第三类场景以冷热共用设备、储热、建筑RC与配电边界建立原型服务模型。其规模、耗时、gap和审计残差属于软件验证；未经来源准入的容量、费用及排放属于情景数值，不形成真实工程改善结论。

主证据包括NASA产品全量来源复核、N2首轮与探索性比较，以及原作者多目标例题复现。投资运行的软件记录仍按代表规划、完整冷热服务年和供暖季分别报告。未来工程案例须关联通过来源审查的配置、规划与回放manifest及同条件基准。

\begin{table}[H]
\centering
\caption{最终案例所需的结果证据}
\begin{tabularx}{\textwidth}{@{}p{0.19\textwidth}p{0.28\textwidth}X@{}}
\toprule
记录 & 结果内容 & 核验条件\\
\midrule
规划 & 投资、代表期范围、状态与gap & 代码及来源版本确定，覆盖权重和建筑边界可追溯\\
年度 & 完整冷热服务、费用、总碳与地源 & 连续逐时、投资冻结、资本一次计入、逐源年度账本\\
供暖季 & 该期间交付热量、状态与运行费 & 独立时钟和边界；不与年度资本或全年地源条件混用\\
建筑审计 & 温度、RC残差、冷热互斥及交付量 & 轨迹同源、时段末采样范围明确；观测精度另验证\\
公平基准 & 同服务条件下的差值 & 同天气、建筑、舒适、端点和成本碳口径，来源准入通过\\
\bottomrule
\end{tabularx}
\end{table}

\subsection{公平比较与外部证据}

基准和推荐方案采用相同气象、建筑参数、服务需求、舒适与初末状态、失供要求及会计边界。技术限制或容量调整必须预先声明，不能通过降低服务质量、删除冷需求或移除资本使某方案显得更优。基准未完成同条件核验时，仅展示程序状态，不给出改善结论。

外部RC观测的留出验证与投资调度比较分开报告。只有真实观测、参数辨识划分和留出预测全部可追溯时才计算外部误差；无法获得观测的参数继续保留假设身份。模型求解记录及程序审计不作为企业认可、鉴定或现场测量的替代材料。

\Needspace{8\baselineskip}
\subsection{独立外部RC适用性验证}

采用Fraunhofer/IEA EBC Annex 71德国N2主实验发布数据\supcite{8}，对单一1R1C结构作外部适用性诊断。数据版本2、DOI与CC BY-SA 4.0许可可查，全部11988条连续记录按时间分为训练、验证和测试段，数量分别为7192、2397和2399。独立复算核对来源文件及162项指标，首轮结果已获有限方法引用准入。

该资料是发布方处理后的测量数据。短缺失温度插值及部分功率、流量置零没有逐行质量旗标，无法剔除全部受处理样本，也不能排除源处理跨划分信息。选择N2排除其他房屋明确的模拟补热，不等于获得未经处理的原始传感器数据。

十个房间110~cm测点温度按地板面积加权形成派生观测，不称全屋体积均温。加热及内部热源功率由W转换为kW，电热按1倍换算为显热是声明假设。验证采用区间平均观测方程；室内与功率十分钟聚合窗口、标签相位及计量不确定度尚未完整确认，方程代数精确不证明秒级过程精确。太阳得热采用声明代理，窗面积与SHGC乘积不作为认证全屋严格热增益上界。

训练、验证和测试按连续时间划分，以独立数据检验预测表现\supcite{3}。本次参数与起点选择仅依据训练段；内部方案文件不单独证明不可变事前注册。一步预测使用上一实测区间温度，连续开环仅初始化一次，分段开环仅在各划分段起点初始化；持续温度基准采用同样的信息和初始化规则。

\begin{table}[H]
\centering
\caption{N2首轮与24~h损失探索的测试段RMSE（$^\circ$C）}
\begin{tabularx}{\textwidth}{@{}Xrrr@{}}
\toprule
预测模式与初始化 & 首轮RC & 后续探索RC & 持续性基准\\
\midrule
一步：上一实测区间温度 & 0.085979 & 0.091334 & 0.091964\\
连续开环：全序列初始化一次 & 1.885259 & 7.425830 & 0.790172\\
分段开环：各划分段初始化一次 & 1.155136 & 6.170260 & 0.626294\\
\bottomrule
\end{tabularx}
\end{table}

首轮一步误差略低于基准，两种开环均高于基准，揭示长时漂移。后续采用24~h日块训练损失的探索沿用同一划分，测试开环进一步恶化，太阳代理参数达到上界；训练改善未转化为留出改善，不能称新独立验证。完整重拟合独立复现首轮与后续参数及指标。结果仅支持已知天气和热输入下有限的条件单步适用性；德国拟合H/C不迁移到武汉，也不证明中国工程效果。

\clearpage
\subsection{原作者多目标算例复现}

采用Mavrotas作者公开报告的双目标连续LP例题\supcite{5}：
\begin{equation}
\max f_1=x_1,\qquad \max f_2=3x_1+4x_2,
\end{equation}
\begin{equation}
0\leq x_1\leq20,\quad0\leq x_2\leq40,\quad5x_1+4x_2\leq200.
\end{equation}
非负域依据原图2解释，其余系数按原式录入。复现在项目LinearMILP数值层上另用两阶段包装：先最大化主目标，再等式固定实际所得主目标值并优化次目标。两端点$(f_1,f_2)$为$(20,160)$与$(8,184)$，与作者表2一致。

由实际端点的$f_2$范围四等分，逐档施加$f_2\geq\varepsilon$并分层求解。表\ref{tab:published-epsilon}给出实际结果，五点唯一且互不支配，与独立有理数多边形解析界一致。
\begin{table}[H]
\centering
\caption{公开数学例题的实际分层$\varepsilon$约束解}
\label{tab:published-epsilon}
\begin{tabular}{rrrrr}
\toprule
$\varepsilon$ & $x_1$ & $x_2$ & $f_1$ & $f_2$\\
\midrule
160 & 20 & 25 & 20 & 160\\
166 & 17 & 28.75 & 17 & 166\\
172 & 14 & 32.5 & 14 & 172\\
178 & 11 & 36.25 & 11 & 178\\
184 & 8 & 40 & 8 & 184\\
\bottomrule
\end{tabular}
\end{table}

\begin{figure}[H]
\centering
\begin{tikzpicture}[x=0.5cm,y=0.13cm,font=\small]
\draw[arrow] (-1,-2) -- (14,-2) node[right] {$f_1$（最大化）};
\draw[arrow] (-1,-2) -- (-1,28) node[above] {$f_2$（最大化）};
\foreach \x/\lab in {0/8,3/11,6/14,9/17,12/20}{
  \draw (\x,-2) -- (\x,-2.7) node[below] {\lab};
}
\foreach \y/\lab in {0/160,6/166,12/172,18/178,24/184}{
  \draw (-1,\y) -- (-1.2,\y) node[left] {\lab};
}
\draw[thick,draw=paperGreenDark] (0,24) -- (12,0);
\foreach \x/\y in {0/24,3/18,6/12,9/6,12/0}{
  \fill[paperGreenDark] (\x,\y) circle[radius=2pt];
}
\node[align=left] at (8,23) {实点：实际LP解\\实线：独立解析前沿};
\end{tikzpicture}
\caption{公开数学例题的五个采样点与解析前沿}
\end{figure}

含资源上限180/220的声明边界检验，共实际调用23次LP：22次最优，1次预期不可行；最大矩阵行、变量界违反及目标重算误差均为0。累计求解0.017743~s属于此小例题，不代表工程求解速度；未暴露的LP gap保留null。复现未调用\path{HeatPlanningModel.solve_lexicographic}，不属于完整AUGMECON或供热工程收益验证。项目最终全量回归628项通过，耗时67.55~s，只作为软件检查记录。
